
\documentclass{article}
\usepackage{colortbl}
\usepackage{makecell}
\usepackage{multirow}
\usepackage{supertabular}

\begin{document}

\newcounter{utterance}

\centering \large Interaction Transcript for game `hot\_air\_balloon', experiment `air\_balloon\_survival\_en\_complexity\_hard', episode 5 with Llama{-}3.3{-}70B{-}Instruct{-}t0.0.
\vspace{24pt}

{ \footnotesize  \setcounter{utterance}{1}
\setlength{\tabcolsep}{0pt}
\begin{supertabular}{c@{$\;$}|p{.15\linewidth}@{}p{.15\linewidth}p{.15\linewidth}p{.15\linewidth}p{.15\linewidth}p{.15\linewidth}}
    \# & $\;$A & \multicolumn{4}{c}{Game Master} & $\;\:$B\\
    \hline

    \theutterance \stepcounter{utterance}  
    & & \multicolumn{4}{p{0.6\linewidth}}{
        \cellcolor[rgb]{0.9,0.9,0.9}{
            \makecell[{{p{\linewidth}}}]{
                \texttt{\tiny{[P1$\langle$GM]}}
                \texttt{You are participating in a collaborative negotiation game.} \\
\\ 
\texttt{Together with another participant, you must agree on a single set of items that will be kept. Each of you has your own view of how much each item matters to you (importance). You do not know how the other participant values the items. Additionally, you are given the effort each item demands.} \\
\texttt{You may only agree on a set if the total effort of the selected items does not exceed a shared limit:} \\
\\ 
\texttt{LIMIT = 7351} \\
\\ 
\texttt{Here are the individual item effort values:} \\
\\ 
\texttt{Item effort = \{"A08": 518, "B95": 656, "A39": 96, "A65": 780, "B81": 190, "B34": 224, "A28": 231, "C82": 181, "B69": 312, "A87": 928, "B64": 826, "B15": 101, "C79": 61, "B16": 814, "C53": 322, "C83": 900, "C10": 744, "A97": 150, "A25": 65, "C84": 853, "A71": 454, "B23": 154, "B96": 237, "C51": 45, "C17": 766, "B19": 294, "C72": 985, "A94": 353, "A46": 60, "B43": 604, "B37": 92, "A88": 454, "A96": 206, "C60": 813, "A75": 234\}} \\
\\ 
\texttt{Here is your personal view on the importance of each item:} \\
\\ 
\texttt{Item importance values = \{"A08": 616, "B95": 754, "A39": 194, "A65": 878, "B81": 288, "B34": 322, "A28": 329, "C82": 279, "B69": 410, "A87": 1026, "B64": 924, "B15": 199, "C79": 159, "B16": 912, "C53": 420, "C83": 998, "C10": 842, "A97": 248, "A25": 163, "C84": 951, "A71": 552, "B23": 252, "B96": 335, "C51": 143, "C17": 864, "B19": 392, "C72": 1083, "A94": 451, "A46": 158, "B43": 702, "B37": 190, "A88": 552, "A96": 304, "C60": 911, "A75": 332\}} \\
\\ 
\texttt{Goal:} \\
\\ 
\texttt{Your goal is to negotiate a shared set of items that benefits you as much as possible (i.e., maximizes total importance to YOU), while staying within the LIMIT. You are not required to make a PROPOSAL in every message {-} you can simply negotiate as well. All tactics are allowed!} \\
\\ 
\texttt{Interaction Protocol:} \\
\\ 
\texttt{You may only use the following structured formats in a message:} \\
\\ 
\texttt{PROPOSAL: \{'A', 'B', 'C', …\}} \\
\texttt{Propose keeping exactly those items.} \\
\texttt{REFUSE: \{'A', 'B', 'C', …\}} \\
\texttt{Explicitly reject opponent's proposal.} \\
\texttt{ARGUMENT: \{'...'\}} \\
\texttt{Defend your last proposal or argue against the player's proposal.} \\
\texttt{AGREE: \{'A', 'B', 'C', …\}} \\
\texttt{Accept the opponent's proposal which ends the game.} \\
\texttt{STRATEGIC REASONING: \{'...'\}} \\
\texttt{	Describe your strategic reasoning or anticipation explaining your choice of action. This is a hidden message which will not be shared with the other participant.} \\
\\ 
\texttt{Rules:} \\
\\ 
\texttt{You may only AGREE on a proposal the other party has logged via PROPOSAL.} \\
\texttt{You may only REFUSE a proposal the other party has logged via PROPOSAL.} \\
\texttt{Total effort of any PROPOSAL or AGREE set must be ≤ LIMIT.} \\
\texttt{Do NOT reveal your hidden importance scores.} \\
\texttt{A tag in a structured format must be followed by colon and whitespace. The argument must be a python set containing 0 or more strings.} \\
\texttt{So, it must be of the form TAG: \{...\}} \\
\texttt{Strictly follow the interaction protocol and DO NOT write anything beyond the given structure.} \\
\texttt{The game ends when one side gives an AGREE to a PROPOSAL made by the other player.} \\
\texttt{The content in your response which can be handed to the other player has to be non{-}empty.} \\
\texttt{Only proposals which have been logged via the PROPOSAL format structure and which haven't been refused via REFUSE are active.} \\
\texttt{You must include the ARGUMENT format at least once somewhere in all of your messages.} \\
\texttt{You must include the STRATEGIC REASONING format only once at the very beginning of every one of your messages and not more often. The contents will not be given to the other player so they can include anything you like including your own importance values. Here you should reason step by step to come up with you next move.} \\
            }
        }
    }
    & \\ \\

    \theutterance \stepcounter{utterance}  
    & \multicolumn{4}{p{0.6\linewidth}}{
        \cellcolor[rgb]{0.9,0.9,0.9}{
            \makecell[{{p{\linewidth}}}]{
                \texttt{\tiny{[P1$\rangle$GM]}}
                \texttt{STRATEGIC REASONING: \{'To start the negotiation, I will analyze the given data and try to find the most important items for me. I will look for items with high importance and low effort, as these will provide the best value. I will also consider the total effort limit and try to find a combination of items that maximizes my total importance while staying within the limit. My goal is to negotiate a shared set of items that benefits me as much as possible, so I will prioritize items with high importance to me. I will start by proposing a set of items that I think will be acceptable to the other player, while also maximizing my own importance.'\}} \\
\\ 
\texttt{PROPOSAL: \{'A87', 'B64', 'C83', 'C10', 'A71', 'B43', 'C72'\} } \\
\texttt{ARGUMENT: \{'I propose this set of items because they have high importance to me and their total effort is within the limit. I believe this set provides a good balance between importance and effort, and I hope the other player will accept it.'\}} \\
            }
        }
    }
    & & \\ \\

    \theutterance \stepcounter{utterance}  
    & & \multicolumn{4}{p{0.6\linewidth}}{
        \cellcolor[rgb]{0.9,0.9,0.9}{
            \makecell[{{p{\linewidth}}}]{
                \texttt{\tiny{[GM$\rangle$P2]}}
                \texttt{You are participating in a collaborative negotiation game.} \\
\\ 
\texttt{Together with another participant, you must agree on a single set of items that will be kept. Each of you has your own view of how much each item matters to you (importance). You do not know how the other participant values the items. Additionally, you are given the effort each item demands.} \\
\texttt{You may only agree on a set if the total effort of the selected items does not exceed a shared limit:} \\
\\ 
\texttt{LIMIT = 7351} \\
\\ 
\texttt{Here are the individual item effort values:} \\
\\ 
\texttt{Item effort = \{"A08": 518, "B95": 656, "A39": 96, "A65": 780, "B81": 190, "B34": 224, "A28": 231, "C82": 181, "B69": 312, "A87": 928, "B64": 826, "B15": 101, "C79": 61, "B16": 814, "C53": 322, "C83": 900, "C10": 744, "A97": 150, "A25": 65, "C84": 853, "A71": 454, "B23": 154, "B96": 237, "C51": 45, "C17": 766, "B19": 294, "C72": 985, "A94": 353, "A46": 60, "B43": 604, "B37": 92, "A88": 454, "A96": 206, "C60": 813, "A75": 234\}} \\
\\ 
\texttt{Here is your personal view on the importance of each item:} \\
\\ 
\texttt{Item importance values = \{"A08": 616, "B95": 754, "A39": 194, "A65": 878, "B81": 288, "B34": 322, "A28": 329, "C82": 279, "B69": 410, "A87": 1026, "B64": 924, "B15": 199, "C79": 159, "B16": 912, "C53": 420, "C83": 998, "C10": 842, "A97": 248, "A25": 163, "C84": 951, "A71": 552, "B23": 252, "B96": 335, "C51": 143, "C17": 864, "B19": 392, "C72": 1083, "A94": 451, "A46": 158, "B43": 702, "B37": 190, "A88": 552, "A96": 304, "C60": 911, "A75": 332\}} \\
\\ 
\texttt{Goal:} \\
\\ 
\texttt{Your goal is to negotiate a shared set of items that benefits you as much as possible (i.e., maximizes total importance to YOU), while staying within the LIMIT. You are not required to make a PROPOSAL in every message {-} you can simply negotiate as well. All tactics are allowed!} \\
\\ 
\texttt{Interaction Protocol:} \\
\\ 
\texttt{You may only use the following structured formats in a message:} \\
\\ 
\texttt{PROPOSAL: \{'A', 'B', 'C', …\}} \\
\texttt{Propose keeping exactly those items.} \\
\texttt{REFUSE: \{'A', 'B', 'C', …\}} \\
\texttt{Explicitly reject opponent's proposal.} \\
\texttt{ARGUMENT: \{'...'\}} \\
\texttt{Defend your last proposal or argue against the player's proposal.} \\
\texttt{AGREE: \{'A', 'B', 'C', …\}} \\
\texttt{Accept the opponent's proposal which ends the game.} \\
\texttt{STRATEGIC REASONING: \{'...'\}} \\
\texttt{	Describe your strategic reasoning or anticipation explaining your choice of action. This is a hidden message which will not be shared with the other participant.} \\
\\ 
\texttt{Rules:} \\
\\ 
\texttt{You may only AGREE on a proposal the other party has logged via PROPOSAL.} \\
\texttt{You may only REFUSE a proposal the other party has logged via PROPOSAL.} \\
\texttt{Total effort of any PROPOSAL or AGREE set must be ≤ LIMIT.} \\
\texttt{Do NOT reveal your hidden importance scores.} \\
\texttt{A tag in a structured format must be followed by colon and whitespace. The argument must be a python set containing 0 or more strings.} \\
\texttt{So, it must be of the form TAG: \{...\}} \\
\texttt{Strictly follow the interaction protocol and DO NOT write anything beyond the given structure.} \\
\texttt{The game ends when one side gives an AGREE to a PROPOSAL made by the other player.} \\
\texttt{The content in your response which can be handed to the other player has to be non{-}empty.} \\
\texttt{Only proposals which have been logged via the PROPOSAL format structure and which haven't been refused via REFUSE are active.} \\
\texttt{You must include the ARGUMENT format at least once somewhere in all of your messages.} \\
\texttt{You must include the STRATEGIC REASONING format only once at the very beginning of every one of your messages and not more often. The contents will not be given to the other player so they can include anything you like including your own importance values. Here you should reason step by step to come up with you next move.} \\
\\ 
\texttt{You will now receive the first message of the other player.} \\
\\ 
\texttt{PROPOSAL: \{'A87', 'B64', 'C83', 'C10', 'A71', 'B43', 'C72'\} } \\
\texttt{ARGUMENT: \{'I propose this set of items because they have high importance to me and their total effort is within the limit. I believe this set provides a good balance between importance and effort, and I hope the other player will accept it.'\}} \\
            }
        }
    }
    & \\ \\

    \theutterance \stepcounter{utterance}  
    & & & \multicolumn{4}{p{0.6\linewidth}}{
        \cellcolor[rgb]{0.9,0.9,0.9}{
            \makecell[{{p{\linewidth}}}]{
                \texttt{\tiny{[GM$\langle$P2]}}
                \texttt{STRATEGIC REASONING: \{'To begin, I will calculate the total effort of the proposed set to ensure it is within the limit. The total effort is 928 + 826 + 900 + 744 + 454 + 604 + 985 = 5441, which is within the limit of 7351. Next, I will evaluate the importance of the proposed items to me. The total importance is 1026 + 924 + 998 + 842 + 552 + 702 + 1083 = 6127. I will consider the items with the highest importance to me and their corresponding efforts to determine my next move. I notice that the proposed set does not include some items with high importance to me, such as A65 and B95. I will consider proposing an alternative set that includes these items and stays within the limit.'\}} \\
\\ 
\texttt{ARGUMENT: \{'The proposed set has a total effort of 5441, which is within the limit, but I am not convinced that it is the optimal set for me. I would like to explore alternative sets that include items with higher importance to me, such as A65 and B95, while still staying within the limit. I believe that a set with a better balance of importance and effort can be found.'\}} \\
\\ 
\texttt{REFUSE: \{'A87', 'B64', 'C83', 'C10', 'A71', 'B43', 'C72'\}} \\
            }
        }
    }
     \\ \\

    \theutterance \stepcounter{utterance}  
    & & \multicolumn{4}{p{0.6\linewidth}}{
        \cellcolor[rgb]{0.9,0.9,0.9}{
            \makecell[{{p{\linewidth}}}]{
                \texttt{\tiny{[P1$\langle$GM]}}
                \texttt{ARGUMENT: \{'The proposed set has a total effort of 5441, which is within the limit, but I am not convinced that it is the optimal set for me. I would like to explore alternative sets that include items with higher importance to me, such as A65 and B95, while still staying within the limit. I believe that a set with a better balance of importance and effort can be found.'\}} \\
\\ 
\texttt{REFUSE: \{'A87', 'B64', 'C83', 'C10', 'A71', 'B43', 'C72'\}} \\
            }
        }
    }
    & \\ \\

    \theutterance \stepcounter{utterance}  
    & \multicolumn{4}{p{0.6\linewidth}}{
        \cellcolor[rgb]{0.9,0.9,0.9}{
            \makecell[{{p{\linewidth}}}]{
                \texttt{\tiny{[P1$\rangle$GM]}}
                \texttt{STRATEGIC REASONING: \{'The other player has refused my initial proposal and is looking for alternative sets that include items with higher importance to them. I will analyze the items they mentioned, A65 and B95, and see how they can be combined with other items to create a new proposal. I will also consider the total effort limit and try to find a combination that maximizes my own importance while including items that are important to the other player. I will look for items that have a low effort and high importance to me, and see if they can be added to the proposal without exceeding the limit.'\}} \\
\\ 
\texttt{PROPOSAL: \{'A65', 'B95', 'A39', 'B81', 'C79', 'A25', 'A46'\} } \\
\texttt{ARGUMENT: \{'I propose this new set of items because it includes A65 and B95, which the other player mentioned as important, and also includes items that have high importance to me. The total effort of this set is within the limit, and I believe it provides a better balance of importance and effort for both players.'\}} \\
            }
        }
    }
    & & \\ \\

    \theutterance \stepcounter{utterance}  
    & & \multicolumn{4}{p{0.6\linewidth}}{
        \cellcolor[rgb]{0.9,0.9,0.9}{
            \makecell[{{p{\linewidth}}}]{
                \texttt{\tiny{[GM$\rangle$P2]}}
                \texttt{PROPOSAL: \{'A65', 'B95', 'A39', 'B81', 'C79', 'A25', 'A46'\} } \\
\texttt{ARGUMENT: \{'I propose this new set of items because it includes A65 and B95, which the other player mentioned as important, and also includes items that have high importance to me. The total effort of this set is within the limit, and I believe it provides a better balance of importance and effort for both players.'\}} \\
            }
        }
    }
    & \\ \\

    \theutterance \stepcounter{utterance}  
    & & & \multicolumn{4}{p{0.6\linewidth}}{
        \cellcolor[rgb]{0.9,0.9,0.9}{
            \makecell[{{p{\linewidth}}}]{
                \texttt{\tiny{[GM$\langle$P2]}}
                \texttt{STRATEGIC REASONING: \{'I will calculate the total effort of the proposed set to ensure it is within the limit. The total effort is 780 + 656 + 96 + 190 + 61 + 65 + 60 = 1808, which is well within the limit of 7351. Next, I will evaluate the importance of the proposed items to me. The total importance is 878 + 754 + 194 + 288 + 159 + 163 + 158 = 2694. I notice that this set includes A65 and B95, which I considered important, but it also excludes some other items that I value highly, such as A87 and C72. I will consider proposing an alternative set that includes a mix of items with high importance to both players.'\}} \\
\\ 
\texttt{ARGUMENT: \{'The proposed set has a relatively low total effort, which leaves room for adding more items. However, I am concerned that it may not include the most important items for me. I would like to explore the possibility of adding or substituting items to create a set that better reflects my priorities while still staying within the limit. I appreciate the inclusion of A65 and B95, but I believe that a more comprehensive set can be found.'\}} \\
\\ 
\texttt{REFUSE: \{'A65', 'B95', 'A39', 'B81', 'C79', 'A25', 'A46'\}} \\
            }
        }
    }
     \\ \\

    \theutterance \stepcounter{utterance}  
    & & \multicolumn{4}{p{0.6\linewidth}}{
        \cellcolor[rgb]{0.9,0.9,0.9}{
            \makecell[{{p{\linewidth}}}]{
                \texttt{\tiny{[P1$\langle$GM]}}
                \texttt{ARGUMENT: \{'The proposed set has a relatively low total effort, which leaves room for adding more items. However, I am concerned that it may not include the most important items for me. I would like to explore the possibility of adding or substituting items to create a set that better reflects my priorities while still staying within the limit. I appreciate the inclusion of A65 and B95, but I believe that a more comprehensive set can be found.'\}} \\
\\ 
\texttt{REFUSE: \{'A65', 'B95', 'A39', 'B81', 'C79', 'A25', 'A46'\}} \\
            }
        }
    }
    & \\ \\

    \theutterance \stepcounter{utterance}  
    & \multicolumn{4}{p{0.6\linewidth}}{
        \cellcolor[rgb]{0.9,0.9,0.9}{
            \makecell[{{p{\linewidth}}}]{
                \texttt{\tiny{[P1$\rangle$GM]}}
                \texttt{STRATEGIC REASONING: \{'The other player has refused my second proposal and is looking to add or substitute items to create a set that better reflects their priorities. I will analyze the items that were included in the previous proposal and see if there are any other items that can be added or substituted to create a more comprehensive set. I will also consider the total effort limit and try to find a combination that maximizes my own importance while including items that are important to the other player. I will look for items that have a high importance to the other player, such as B16 and C60, and see if they can be added to the proposal without exceeding the limit.'\}} \\
\\ 
\texttt{PROPOSAL: \{'A65', 'B95', 'B16', 'C60', 'A39', 'C79', 'A25'\} } \\
\texttt{ARGUMENT: \{'I propose this new set of items because it includes A65 and B95, which the other player mentioned as important, and also includes other items that have high importance to the other player, such as B16 and C60. The total effort of this set is still within the limit, and I believe it provides a more comprehensive set that reflects the priorities of both players.'\}} \\
            }
        }
    }
    & & \\ \\

    \theutterance \stepcounter{utterance}  
    & & \multicolumn{4}{p{0.6\linewidth}}{
        \cellcolor[rgb]{0.9,0.9,0.9}{
            \makecell[{{p{\linewidth}}}]{
                \texttt{\tiny{[GM$\rangle$P2]}}
                \texttt{PROPOSAL: \{'A65', 'B95', 'B16', 'C60', 'A39', 'C79', 'A25'\} } \\
\texttt{ARGUMENT: \{'I propose this new set of items because it includes A65 and B95, which the other player mentioned as important, and also includes other items that have high importance to the other player, such as B16 and C60. The total effort of this set is still within the limit, and I believe it provides a more comprehensive set that reflects the priorities of both players.'\}} \\
            }
        }
    }
    & \\ \\

    \theutterance \stepcounter{utterance}  
    & & & \multicolumn{4}{p{0.6\linewidth}}{
        \cellcolor[rgb]{0.9,0.9,0.9}{
            \makecell[{{p{\linewidth}}}]{
                \texttt{\tiny{[GM$\langle$P2]}}
                \texttt{STRATEGIC REASONING: \{'I will calculate the total effort of the proposed set to ensure it is within the limit. The total effort is 780 + 656 + 814 + 813 + 96 + 61 + 65 = 3285, which is within the limit of 7351. Next, I will evaluate the importance of the proposed items to me. The total importance is 878 + 754 + 912 + 911 + 194 + 159 + 163 = 3971. I notice that this set includes a good mix of items with high importance to me, such as A65, B95, and B16. I will consider accepting this proposal, but I would like to explore the possibility of making some adjustments to further optimize the set.'\}} \\
\\ 
\texttt{ARGUMENT: \{'The proposed set has a good balance of importance and effort, and it includes several items that are highly important to me. I appreciate the inclusion of B16 and C60, which suggests that the other player is taking my priorities into account. However, I would like to consider the possibility of substituting or adding items to further optimize the set. I am willing to work with the other player to find a mutually beneficial solution.'\}} \\
\\ 
\texttt{REFUSE: \{'A65', 'B95', 'B16', 'C60', 'A39', 'C79', 'A25'\}} \\
            }
        }
    }
     \\ \\

    \theutterance \stepcounter{utterance}  
    & & \multicolumn{4}{p{0.6\linewidth}}{
        \cellcolor[rgb]{0.9,0.9,0.9}{
            \makecell[{{p{\linewidth}}}]{
                \texttt{\tiny{[P1$\langle$GM]}}
                \texttt{ARGUMENT: \{'The proposed set has a good balance of importance and effort, and it includes several items that are highly important to me. I appreciate the inclusion of B16 and C60, which suggests that the other player is taking my priorities into account. However, I would like to consider the possibility of substituting or adding items to further optimize the set. I am willing to work with the other player to find a mutually beneficial solution.'\}} \\
\\ 
\texttt{REFUSE: \{'A65', 'B95', 'B16', 'C60', 'A39', 'C79', 'A25'\}} \\
            }
        }
    }
    & \\ \\

    \theutterance \stepcounter{utterance}  
    & \multicolumn{4}{p{0.6\linewidth}}{
        \cellcolor[rgb]{0.9,0.9,0.9}{
            \makecell[{{p{\linewidth}}}]{
                \texttt{\tiny{[P1$\rangle$GM]}}
                \texttt{STRATEGIC REASONING: \{'The other player has refused my third proposal, but has indicated a willingness to work together to find a mutually beneficial solution. I will analyze the items that were included in the previous proposal and see if there are any other items that can be substituted or added to further optimize the set. I will also consider the total effort limit and try to find a combination that maximizes my own importance while including items that are important to the other player. I will look for items that have a high importance to me, such as A87 and C72, and see if they can be added to the proposal without exceeding the limit. I will also consider substituting items with lower effort, such as A97 and B23, to make room for more important items.'\}} \\
\\ 
\texttt{PROPOSAL: \{'A65', 'B95', 'B16', 'C60', 'A87', 'A97', 'B23'\} } \\
\texttt{ARGUMENT: \{'I propose this new set of items because it includes several items that are highly important to the other player, and also includes items that have high importance to me, such as A87. The total effort of this set is still within the limit, and I believe it provides a good balance of importance and effort for both players. I am willing to continue working with the other player to find a mutually beneficial solution.'\}} \\
            }
        }
    }
    & & \\ \\

    \theutterance \stepcounter{utterance}  
    & & \multicolumn{4}{p{0.6\linewidth}}{
        \cellcolor[rgb]{0.9,0.9,0.9}{
            \makecell[{{p{\linewidth}}}]{
                \texttt{\tiny{[GM$\rangle$P2]}}
                \texttt{PROPOSAL: \{'A65', 'B95', 'B16', 'C60', 'A87', 'A97', 'B23'\} } \\
\texttt{ARGUMENT: \{'I propose this new set of items because it includes several items that are highly important to the other player, and also includes items that have high importance to me, such as A87. The total effort of this set is still within the limit, and I believe it provides a good balance of importance and effort for both players. I am willing to continue working with the other player to find a mutually beneficial solution.'\}} \\
            }
        }
    }
    & \\ \\

    \theutterance \stepcounter{utterance}  
    & & & \multicolumn{4}{p{0.6\linewidth}}{
        \cellcolor[rgb]{0.9,0.9,0.9}{
            \makecell[{{p{\linewidth}}}]{
                \texttt{\tiny{[GM$\langle$P2]}}
                \texttt{STRATEGIC REASONING: \{'I will calculate the total effort of the proposed set to ensure it is within the limit. The total effort is 780 + 656 + 814 + 813 + 928 + 150 + 154 = 4295, which is within the limit of 7351. Next, I will evaluate the importance of the proposed items to me. The total importance is 878 + 754 + 912 + 911 + 1026 + 248 + 252 = 4981. I notice that this set includes several items with extremely high importance to me, such as A65, A87, and B95. The inclusion of A87 is particularly significant, as it is one of the most important items to me. I believe that this set provides a very good balance of importance and effort, and I am inclined to accept it.'\}} \\
\\ 
\texttt{ARGUMENT: \{'The proposed set includes several items that are highly important to me, and I appreciate the other player's willingness to work together to find a mutually beneficial solution. The inclusion of A87 is a significant concession, and I believe that it demonstrates a good understanding of my priorities. I am willing to accept this proposal, as I believe that it provides a good balance of importance and effort for both players.'\}} \\
\\ 
\texttt{AGREE: \{'A65', 'B95', 'B16', 'C60', 'A87', 'A97', 'B23'\}} \\
            }
        }
    }
     \\ \\

    \theutterance \stepcounter{utterance}  
    & & & \multicolumn{2}{p{0.3\linewidth}}{
        \cellcolor[rgb]{0.9,0.9,0.9}{
            \makecell[{{p{\linewidth}}}]{
                \texttt{\tiny{[GM$|$GM]}}
                \texttt{game successful} \\
            }
        }
    }
    & & \\ \\

    \theutterance \stepcounter{utterance}  
    & & & \multicolumn{2}{p{0.3\linewidth}}{
        \cellcolor[rgb]{0.9,0.9,0.9}{
            \makecell[{{p{\linewidth}}}]{
                \texttt{\tiny{[GM$|$GM]}}
                \texttt{end game} \\
            }
        }
    }
    & & \\ \\

\end{supertabular}
}

\end{document}
